\providecommand{\Res}[2]{\qty[round-mode=figures,round-precision=3]{\fpeval{#1}}{#2}}
\providecommand{\Val}[1]{\num[round-mode=figures,round-precision=4]{\fpeval{#1}}}
\providecommand{\ResSci}[2]{\qty[round-mode=figures,round-precision=3,exponent-mode=scientific]{\fpeval{#1}}{#2}}
\providecommand{\Valt}[1]{\expanded{\noexpand\num[round-mode=figures,round-precision=4,exponent-mode=threshold,exponent-thresholds=-3\string:5]}{\fpeval{#1}}}
\providecommand{\NumT}[1]{\expanded{\noexpand\num[exponent-mode=threshold,exponent-thresholds=-3\string:9]}{#1}}
\providecommand{\QtyT}[2]{\expanded{\noexpand\qty[exponent-mode=threshold,exponent-thresholds=-3\string:9]}{#1}{#2}}
\subsubsection*{Question a) Lois de Kirchhoff}
\begin{itemize}[label= -]
    \item Loi des nœuds en $A$ : $I_1 + I_3 = I_2$
    \item Maille de gauche (sens horaire, à partir de $B$) : $\mathcal{E}_1 - R_1 I_1 - R_2 I_2 = 0$
    \item Maille de droite (sens antihoraire, à partir de $B$) : $\mathcal{E}_2 - R_3 I_3 - R_2 I_2 = 0$
\end{itemize}

\subsubsection*{Question b) Courant $I_2$}
\begin{enumerate}[label=\arabic*)]
    \item \textbf{Variables connues et recherchées :} $\mathcal{E}_1 = \qty{\v{e1}}{V}$, $\mathcal{E}_2 = \qty{\v{e2}}{V}$, $R_1 = \qty{\v{R1}}{\ohm}$, $R_2 = \qty{\v{R2}}{\ohm}$, $R_3 = \qty{\v{R3}}{\ohm}$ ; variable recherchée : $I_2$
    \item \textbf{Équations pertinentes :} les trois équations de la question a).
    \item \textbf{Isolation de la variable recherchée :} des deux mailles, $I_1 = \dfrac{\mathcal{E}_1 - R_2 I_2}{R_1}$ et $I_3 = \dfrac{\mathcal{E}_2 - R_2 I_2}{R_3}$. Dans la loi des nœuds :
    \[ \frac{\mathcal{E}_1 - R_2 I_2}{R_1} + \frac{\mathcal{E}_2 - R_2 I_2}{R_3} = I_2 \quad\Rightarrow\quad I_2 = \frac{\dfrac{\mathcal{E}_1}{R_1} + \dfrac{\mathcal{E}_2}{R_3}}{1 + \dfrac{R_2}{R_1} + \dfrac{R_2}{R_3}} \]
    \item \textbf{Résolution numérique :}
    \[ I_2 = \frac{\dfrac{\num{\v{e1}}}{\num{\v{R1}}} + \dfrac{\num{\v{e2}}}{\num{\v{R3}}}}{1 + \dfrac{\num{\v{R2}}}{\num{\v{R1}}} + \dfrac{\num{\v{R2}}}{\num{\v{R3}}}} = \Res{(\v{e1}/\v{R1}+\v{e2}/\v{R3})/(1+\v{R2}/\v{R1}+\v{R2}/\v{R3})}{A} \]
    \item \textbf{Réponse :} $I_2 = \mathbf{\Res{(\v{e1}/\v{R1}+\v{e2}/\v{R3})/(1+\v{R2}/\v{R1}+\v{R2}/\v{R3})}{A}}$.
\end{enumerate}

\subsubsection*{Question c) Courants $I_1$ et $I_3$}
\begin{enumerate}[label=\arabic*)]
    \item \textbf{Variables connues et recherchées :} $I_2 = \Res{(\v{e1}/\v{R1}+\v{e2}/\v{R3})/(1+\v{R2}/\v{R1}+\v{R2}/\v{R3})}{A}$ ; variables recherchées : $I_1$ et $I_3$
    \item \textbf{Équations pertinentes :} les lois des mailles de la question a).
    \item \textbf{Isolation de la variable recherchée :}
    \[ I_1 = \frac{\mathcal{E}_1 - R_2 I_2}{R_1} \qquad I_3 = \frac{\mathcal{E}_2 - R_2 I_2}{R_3} \]
    \item \textbf{Résolution numérique :}
    \[ I_1 = \frac{\num{\v{e1}} - \num{\v{R2}} \cdot \Valt{(\v{e1}/\v{R1}+\v{e2}/\v{R3})/(1+\v{R2}/\v{R1}+\v{R2}/\v{R3})}}{\num{\v{R1}}} = \Res{(\v{e1}-\v{R2}*(\v{e1}/\v{R1}+\v{e2}/\v{R3})/(1+\v{R2}/\v{R1}+\v{R2}/\v{R3}))/\v{R1}}{A} \]
    \[ I_3 = \frac{\num{\v{e2}} - \num{\v{R2}} \cdot \Valt{(\v{e1}/\v{R1}+\v{e2}/\v{R3})/(1+\v{R2}/\v{R1}+\v{R2}/\v{R3})}}{\num{\v{R3}}} = \Res{(\v{e2}-\v{R2}*(\v{e1}/\v{R1}+\v{e2}/\v{R3})/(1+\v{R2}/\v{R1}+\v{R2}/\v{R3}))/\v{R3}}{A} \]
    \item \textbf{Réponse :} $I_1 = \mathbf{\Res{(\v{e1}-\v{R2}*(\v{e1}/\v{R1}+\v{e2}/\v{R3})/(1+\v{R2}/\v{R1}+\v{R2}/\v{R3}))/\v{R1}}{A}}$ et $I_3 = \mathbf{\Res{(\v{e2}-\v{R2}*(\v{e1}/\v{R1}+\v{e2}/\v{R3})/(1+\v{R2}/\v{R1}+\v{R2}/\v{R3}))/\v{R3}}{A}}$ ; les deux sont positifs, donc dans le sens indiqué. Vérification : $I_1 + I_3 = \Res{(\v{e1}-\v{R2}*(\v{e1}/\v{R1}+\v{e2}/\v{R3})/(1+\v{R2}/\v{R1}+\v{R2}/\v{R3}))/\v{R1}+(\v{e2}-\v{R2}*(\v{e1}/\v{R1}+\v{e2}/\v{R3})/(1+\v{R2}/\v{R1}+\v{R2}/\v{R3}))/\v{R3}}{A} = I_2$.
\end{enumerate}

\subsubsection*{Question d) Puissance fournie par la pile 1}
\begin{enumerate}[label=\arabic*)]
    \item \textbf{Variables connues et recherchées :} $\mathcal{E}_1 = \qty{\v{e1}}{V}$, $I_1 = \Res{(\v{e1}-\v{R2}*(\v{e1}/\v{R1}+\v{e2}/\v{R3})/(1+\v{R2}/\v{R1}+\v{R2}/\v{R3}))/\v{R1}}{A}$ ; variable recherchée : $P$
    \item \textbf{Équations pertinentes :}
    \[ P = \mathcal{E} I \]
    \item \textbf{Isolation de la variable recherchée :} l'équation est déjà isolée.
    \item \textbf{Résolution numérique :}
    \[ P = \num{\v{e1}} \cdot \Valt{(\v{e1}-\v{R2}*(\v{e1}/\v{R1}+\v{e2}/\v{R3})/(1+\v{R2}/\v{R1}+\v{R2}/\v{R3}))/\v{R1}} = \Res{\v{e1}*(\v{e1}-\v{R2}*(\v{e1}/\v{R1}+\v{e2}/\v{R3})/(1+\v{R2}/\v{R1}+\v{R2}/\v{R3}))/\v{R1}}{W} \]
    \item \textbf{Réponse :} la pile 1 fournit $P = \mathbf{\Res{\v{e1}*(\v{e1}-\v{R2}*(\v{e1}/\v{R1}+\v{e2}/\v{R3})/(1+\v{R2}/\v{R1}+\v{R2}/\v{R3}))/\v{R1}}{W}}$.
\end{enumerate}
